\documentclass[11pt]{article}

\usepackage[margin=1in]{geometry}
\usepackage[T1]{fontenc}
\usepackage[utf8]{inputenc}
\usepackage{lmodern}
\usepackage{microtype}
\usepackage{amsmath,amssymb,amsthm}
\usepackage{booktabs}
\usepackage{enumitem}
\usepackage{algorithm}
\usepackage{algpseudocode}
\usepackage{float}
\usepackage{natbib}
\usepackage{comment}
\usepackage{xcolor}
\usepackage[hidelinks]{hyperref}

\newcommand{\thetab}{\theta_0}
\newcommand{\taskvec}{\tau}
\newcommand{\loss}{\mathcal{L}}
\newcommand{\tasks}{\mathcal{T}}
\newcommand{\data}{\mathcal{D}}
\newcommand{\R}{\mathbb{R}}
\newcommand{\ind}{\mathbf{1}}
\DeclareMathOperator{\clip}{clip}
\DeclareMathOperator{\rms}{rms}

\usepackage{amsthm}

\newtheorem{theorem}{Theorem}
\newtheorem{corollary}{Corollary}
\newtheorem{proposition}{Proposition}

\title{Attribution-Patching-Guided Replay Refinement for Task-Vector Merging}
\author{Anonymous Authors \\
\small Draft proposal}
\date{\today}

\begin{document}
\maketitle

\begin{abstract}
We study the task of merging fine-tuned variants of a shared pretrained model when a small labeled replay buffer for each task is available.
After an initial task-arithmetic merge, we refine the merged model with task gradients while using the known task experts as coordinate-wise anchors.
As a parameter-space analogue of attribution patching, the effect of moving the current model toward a task expert is approximated by a first-order loss attribution, and only coordinates for which the expert direction is locally loss-decreasing are used. 

ordinary replay gradient descent is that each coordinate's step is weighted by its distance to the expert, gated to keep only steps that move toward the expert, and clipped so it cannot overshoot.
\end{abstract}

\section{Related Work and Background}

\paragraph{Model merging and weight-space averaging.}
Model merging combines multiple trained or fine-tuned models into a single model in weight space, avoiding the inference-time overhead of ensembling. Early and influential forms include direct parameter averaging, model soups, and Fisher-weighted merging. Model soups average multiple fine-tuned checkpoints and can improve accuracy without increasing inference cost when the checkpoints lie in a compatible basin \citep{wortsman2022modelsoups}. Fisher merging instead weights parameters using Fisher information, motivated by a Laplace approximation to the local posterior around each model \citep{matena2022fisher}. These approaches motivate the broader view that homologous fine-tuned models can often be combined directly in parameter space, but they do not by themselves resolve interference when the merged models specialize in different downstream tasks.

\paragraph{Task vectors.}
Task arithmetic represents the effect of fine-tuning as a task vector
\begin{equation}
    \taskvec_t = \theta_t - \thetab,
\end{equation}
where $\thetab$ is a shared pretrained model and $\theta_t$ is the model fine-tuned on task $t$. A merged model is then constructed as
\begin{equation}
    \theta_{\mathrm{merge}} = \thetab + \sum_{t \in \tasks} \lambda_t \taskvec_t,
\end{equation}
where $\lambda_t$ are scalar or tensor-wise coefficients. Task Arithmetic demonstrated that such vectors can be added or negated to steer model behavior \citep{ilharco2022taskarithmetic}. This formulation is attractive because it requires only checkpoints, but naive addition can over-count redundant directions and introduce sign or magnitude conflicts across tasks.

\paragraph{Interference-aware task-vector merging.}
A major line of work addresses interference between task vectors. TIES-Merging trims small-magnitude updates, resolves sign disagreement, and merges only sign-consistent components \citep{yadav2023ties}. DARE randomly drops most delta parameters and rescales the remaining entries, arguing that fine-tuning deltas are highly redundant \citep{yu2023dare}. DELLA generalizes this drop-and-rescale idea with magnitude-based sampling \citep{deep2024della}. Model Breadcrumbs similarly builds sparse masks by removing both negligible and outlier perturbations \citep{davari2023breadcrumbs}. Localize-and-Stitch explicitly localizes sparse task-relevant regions and stitches only those regions back into the pretrained model \citep{he2024localize}. These methods share the view that not all task-vector coordinates should be merged equally. The present proposal does not introduce another pruning rule; instead, it asks whether a small amount of supervised replay can repair an initial merge while staying close to the known task experts.

\paragraph{Data-dependent merging baselines.}
Several merging methods use small amounts of data or model outputs to estimate better combinations. Activated Parameter Locating (APL) estimates task-vector importance through a data-dependent contrast and then uses the result to guide pruning \citep{kong2024apl}. RegMean formulates merging as a regression problem that minimizes prediction differences between the merged model and the candidate models \citep{jin2023regmean}. AdaMerging learns task-wise or layer-wise merge coefficients by minimizing an unsupervised entropy objective on unlabeled samples \citep{yang2024adamerging}. These methods are important baselines because the proposed method also uses data, although it uses labeled replay buffers only for a few post-merge refinement steps rather than for constructing a pre-merge mask or a closed-form fusion rule.

\paragraph{Attribution patching.}
Attribution patching approximates the effect of an intervention by a first-order Taylor expansion, replacing many explicit interventions by gradient-based scores \citep{syed2024attributionpatching,kramar2024atpstar}. In parameter space, the analogous approximation is a directional derivative of the task loss along a chosen parameter displacement. Here the relevant intervention is not a replacement of one expert by another. It is the movement from the current merged model toward the expert for task $i$. This gives a direct bridge between attribution patching and replay-buffer refinement: the attribution score decides which coordinates of the expert direction are locally predicted to decrease the task loss.

\paragraph{Shared and task-specific structure.}
Several recent methods recognize that task vectors contain both shared and exclusive information. Twin-Merging separates shared and exclusive task knowledge and dynamically combines them at inference time \citep{lu2024twin}. Task Vector Bases clusters similar task vectors to compress and reuse shared directions, providing a theoretical and practical framework for reducing redundancy across related tasks \citep{zeng2025taskvectorbases}. FusionBench documents the need for standardized evaluation across model-fusion methods and includes GLUE-style language tasks among its benchmark families \citep{tang2024fusionbench}. Our proposal is complementary: rather than interpreting small attribution values as evidence that a coordinate can be pruned, we use the sign of a merge-to-expert attribution only as a local gate for refinement.

\paragraph{Benchmark settings.}
We use the standard GLUE-8 merging setting as the first benchmark: CoLA, SST-2, MRPC, STS-B, QQP, MNLI, QNLI, and RTE. GLUE was introduced as a multi-task benchmark for natural language understanding \citep{wang2018glue}. We use RoBERTa-base as the primary encoder backbone because it is a strong, widely used model for GLUE-style evaluation \citep{liu2019roberta}. In a full fine-tuning setup, each task has a task-specific classifier or regression head. The shared encoder is merged and evaluated with the corresponding task-specific head. This multi-head setup is convenient and common. %Later experiments should include a shared-output setting, for example GLUE or SuperGLUE in a text-to-text T5-style format \citep{wang2019superglue,raffel2020t5}.

A second-modality benchmark is the standard CLIP/ViT task-vector suite with image-classification tasks such as SUN397, Cars, RESISC45, EuroSAT, SVHN, GTSRB, MNIST, and DTD \citep{radford2021clip,ilharco2022taskarithmetic,tang2024fusionbench}.

\section{Methodology}

\subsection{Problem Setup}
Let $\tasks = \{1,\ldots,T\}$ denote a set of tasks. All task experts are initialized from the same pretrained encoder $\thetab \in \R^d$. Fine-tuning on task $t$ gives expert encoder parameters $\theta_t$ and task vector
\begin{equation}
    \taskvec_t = \theta_t - \thetab.
    \label{eq:task-vector}
\end{equation}
The mergeable encoder parameters are treated as a vector in $\R^d$. We write $r \in \{1,\ldots,d\}$ for a parameter coordinate. Task-specific heads are not merged in the primary GLUE experiments; the merged encoder is evaluated with the corresponding head for each task.

The initial merged encoder is
\begin{equation}
    \theta^{(0)} = \thetab + \sum_{i=1}^{T}\lambda_i\taskvec_i,
    \label{eq:initial-merge}
\end{equation}
where $\lambda_i$ are fixed task-arithmetic coefficients, typically uniform or chosen by a small grid search. Let $\data_t^{\mathrm{probe}}$ be a small replay buffer used for refinement, and let $\data_t^{\mathrm{eval}}$ be the evaluation split used only for reporting performance. For each task $i$, let
\begin{equation}
    \loss_i(\theta;\data_i^{\mathrm{probe}})
    \label{eq:task-loss}
\end{equation}
denote the replay-buffer loss obtained by using encoder parameters $\theta$ with task $i$'s prediction head.
The proposed method does not use a base-relative saliency term and does not construct a pre-merge pruning mask. Its only data-dependent operation is a short replay-buffer refinement of the initial merge.



%The attribution gate that the first derives from a Taylor expansion is exactly the trust-region acceptance rule that the second derives from staying close to the experts, and both yield the identical update in Equations~\eqref{eq:clipped-update}--\eqref{eq:sequential-single-task-update} and hence the same Algorithm~\ref{alg:expert-gated-refinement}.



\subsection{Trust-Region Replay Refinement}
\label{sec:replay-refinement}
Starting from the task-arithmetic initialization in Equation~\eqref{eq:initial-merge}, the natural data-dependent baseline is ordinary replay-buffer gradient descent, which takes an unconstrained step along the task gradients:
\begin{equation}
    \theta^{(s+1)}
    =
    \theta^{(s)}
    -
    \eta \sum_{i=1}^{T}
    \nabla_\theta \loss_i(\theta^{(s)};\data_i^{\mathrm{probe}}).
    \label{eq:naive-replay-gd}
\end{equation}

TODO: List the disadvantages of ordinary GD in this setting here.

The main insight of this paper is that the experts can be used as anchors in this phase.
Rather than taking a free gradient step, every coordinate of the update is tied to the corresponding task expert. For each task $i$ and coordinate $r$ we (i) weight the negative-gradient step by the current distance $|v_{i,r}|$ to the expert, so that coordinates already close to the expert barely move while coordinates far from it move more; (ii) keep the step only where it moves the model toward the expert, discarding coordinates whose replay gradient would pull away from it; and (iii) clip the resulting step to the expert-distance trust region so that a single update can never overshoot the expert. These three operations-are what distinguish the refinement from ordinary replay gradient descent, and together they keep the whole refinement anchored to the known experts.

%A secondary and comparatively minor design choice concerns how the per-task updates are scheduled within a refinement step. An aggregated expert-guided update would compute all $u_i$ at $\theta^{(s)}$ and then apply $\eta\sum_i u_i$ at once; we instead apply each task's update immediately after computing it, so that its gradient, expert direction, and gate are evaluated at the same model state to which it is applied.

We now give more details on our method: Let $\theta^{(s,0)}=\theta^{(s)}$, and let $\pi_s$ be the task order for sweep $s$; $\pi_s$ can be fixed, cyclically rotated, or randomly permuted. For within-sweep index $k\in\{1,\ldots,T\}$, set $i=\pi_s(k)$ and compute
\begin{align}
    g_i^{(s,k)}
    &=
    \nabla_\theta \loss_i(\theta^{(s,k-1)};\data_i^{\mathrm{probe}}),
    \label{eq:sequential-gradient}\\
    v_i^{(s,k)}
    &=
    \theta_i - \theta^{(s,k-1)},
    \label{eq:sequential-direction}\\
    m_{i,r}^{(s,k)}
    &=
    \ind\{g_{i,r}^{(s,k)}v_{i,r}^{(s,k)} < -\epsilon_{\mathrm{gate}}\},
    \label{eq:sequential-mask}\\
    \widetilde u_{i,r}^{(s,k)}
    &=
    -g_{i,r}^{(s,k)}|v_{i,r}^{(s,k)}|m_{i,r}^{(s,k)}.
    \label{eq:sequential-raw-update}
\end{align}
We keep only those coordinates whose replay-gradient step also moves the current model toward expert $i$, and discard the coordinates where the replay step would pull away from the expert and out of the expert-anchored trust region.% This is the same condition $g_{i,r}^{(s,k)}v_{i,r}^{(s,k)}<0$ that Section~\ref{sec:ap-view} obtained from the attribution-patching sign in Equation~\eqref{eq:ap-gate}, now read as \emph{stay within the expert-anchored trust region} instead of \emph{move along a locally loss-decreasing intervention}. The two perspectives therefore agree coordinate by coordinate.
The raw update is first scaled by the learning rate $\eta$ and \emph{then} clipped by the the current distance to the expert, so that the applied per-coordinate step is bounded by $|v_{i,r}^{(s,k)}|$,
\begin{equation}
    u_{i,r}^{(s,k)}
    =
    \clip\!\left(
        \eta\,\widetilde u_{i,r}^{(s,k)},
        - |v_{i,r}^{(s,k)}|,
         |v_{i,r}^{(s,k)}|
    \right),
    \label{eq:clipped-update}
\end{equation}
and the current model is immediately updated with this clipped step:
\begin{equation}
    \theta^{(s,k)}
    =
    \theta^{(s,k-1)} + u_i^{(s,k)}.
    \label{eq:sequential-single-task-update}
\end{equation}
After all tasks in the sweep have been processed, we set
\begin{equation}
    \theta^{(s+1)}=\theta^{(s,T)}.
    \label{eq:sequential-sweep-update}
\end{equation}
%The robustness of the refinement comes primarily from the expert anchoring rather than from this scheduling: the distance weighting and gate keep every coordinate pointed at its expert, and the trust-region clip caps how far it can travel. The immediate application of each update, though secondary to the anchoring, has two motivations of its own. First, each update is applied at the same local model state on which its gradient, expert direction, and gate were computed. Second, the trust-region clipping in Equation~\eqref{eq:clipped-update} controls each single-task movement before the next task sees the modified model, rather than allowing several task vectors to accumulate into one larger coordinate step.

Because the clip is applied \emph{after} the learning-rate scaling, the per-coordinate movement is bounded by $|v_{i,r}^{(s,k)}|$ for any $\eta$; for $\le 1$ a single update can never overshoot the expert in a coordinate, so the refinement stays within the expert-distance trust region and cannot diverge no matter how large $\eta$ is. The learning rate then controls only how many coordinates are driven to the trust-region boundary, not the maximum per-coordinate step size. 

The algorithm is summarized in Algorithm~\ref{alg:expert-gated-refinement}.

%The sequential rule still does not guarantee Pareto improvement across tasks, because an update that helps task $i$ can hurt task $j$. We therefore treat task ordering, cross-task interference, and worst-task retention as explicit diagnostics rather than as consequences of the gate alone. In practice, one can also normalize $u_i^{(s,k)}$ by tensor-wise RMS before applying the single-task update; the defining operation remains the coordinate gate in Equation~\eqref{eq:ap-gate}, whether it is read as an attribution-patching sign or as a trust-region acceptance rule. Under either reading the procedure is the single algorithm summarized in Algorithm~\ref{alg:expert-gated-refinement}.


The following Proposition clarifies the safe region update:

\begin{corollary}[Box-distance monotonicity of the sequential refinement]
\label{cor:algorithm-box-distance}
Let
\[
B_{\mathrm{exp}}
:=
\prod_{r=1}^d
[\underline{\theta}_r,\overline{\theta}_r],
\qquad
\underline{\theta}_r := \min_{j\in\tasks}\theta_{j,r},
\qquad
\overline{\theta}_r := \max_{j\in\tasks}\theta_{j,r},
\]
be the axis-aligned box hull of the task experts. Then every
within-sweep update of Algorithm~\ref{alg:expert-gated-refinement} is
non-expansive in distance to the expert box: for every $1\le p\le\infty$,
\[
    d_p\!\left(\theta^{(s,k)},B_{\mathrm{exp}}\right)
    \le
    d_p\!\left(\theta^{(s,k-1)},B_{\mathrm{exp}}\right).
\]
Consequently, after any number $S$ of sweeps,
\[
    d_p\!\left(\theta^{(S)},B_{\mathrm{exp}}\right)
    \le
    d_p\!\left(\theta^{(0)},B_{\mathrm{exp}}\right)
    \qquad\text{for every }1\le p\le\infty.
\]
\end{corollary}

\begin{proof}
Fix a sweep $s$, a within-sweep index $k$, and let $i=\pi_s(k)$. We first
verify the hypothesis of Theorem~\ref{thm:box-distance-monotonicity} with
\[
    y=\theta^{(s,k-1)},
    \qquad
    z=\theta^{(s,k)},
    \qquad
    x_i=\theta_i.
\]
It is enough to check that, for every coordinate $r$, the new coordinate
$\theta_r^{(s,k)}$ lies between the previous coordinate
$\theta_r^{(s,k-1)}$ and the expert coordinate $\theta_{i,r}$.

Write $v_r=v_{i,r}^{(s,k)}=\theta_{i,r}-\theta_r^{(s,k-1)}$. If the gate is
off, then $m_{i,r}^{(s,k)}=0$, hence $u_{i,r}^{(s,k)}=0$, and the coordinate
is unchanged. If the gate is on, then
\[
    g_{i,r}^{(s,k)} v_r < -\epsilon_{\mathrm{gate}} \le 0,
\]
so $-g_{i,r}^{(s,k)}$ has the same sign as $v_r$. Therefore the raw update
$\widetilde u_{i,r}^{(s,k)}=-g_{i,r}^{(s,k)}|v_r|m_{i,r}^{(s,k)}$ points in
the same direction as $v_r$. Multiplication by $\eta\ge0$ preserves this
direction, possibly making the update zero, and the symmetric clip in
Equation~\eqref{eq:clipped-update} gives an applied update with the same weak sign as $v_r$ and magnitude at most $|v_r|$.
Thus there is an $\alpha_{i,r}^{(s,k)}\in[0,1]$ such that
\[
    u_{i,r}^{(s,k)}=\alpha_{i,r}^{(s,k)}v_{i,r}^{(s,k)}.
\]
The same representation is assumed directly in the optional post-processing
case. Hence
\[
    \theta_r^{(s,k)}
    =
    \theta_r^{(s,k-1)}
    +
    \alpha_{i,r}^{(s,k)}
    \bigl(\theta_{i,r}-\theta_r^{(s,k-1)}\bigr),
\]
so $\theta_r^{(s,k)}$ lies between $\theta_r^{(s,k-1)}$ and $\theta_{i,r}$.

The box $B_{\mathrm{exp}}$ is exactly the axis-aligned hull of the expert
points $\theta_1,\dots,\theta_T$. Therefore
Theorem~\ref{thm:box-distance-monotonicity} gives
\[
    d_p\!\left(\theta^{(s,k)},B_{\mathrm{exp}}\right)
    \le
    d_p\!\left(\theta^{(s,k-1)},B_{\mathrm{exp}}\right)
\]
for every $1\le p\le\infty$. Chaining this inequality over
$k=1,\dots,T$ and over sweeps $s=0,\dots,S-1$ gives the final statement.
\end{proof}



\begin{algorithm}[H]
\caption{Sequential expert-direction-gated replay refinement}
\label{alg:expert-gated-refinement}
\begin{algorithmic}[1]
\Require Pretrained encoder $\thetab$; task experts $\{\theta_i\}_{i=1}^T$; task heads; replay buffers $\{\data_i^{\mathrm{probe}}\}_{i=1}^T$; task-arithmetic weights $\{\lambda_i\}_{i=1}^T$; steps $S$; learning rate $\eta$; gate threshold $\epsilon_{\mathrm{gate}}$.
\State $\taskvec_i \gets \theta_i-\thetab$ for all tasks $i$.
\State $\theta^{(0)} \gets \thetab + \sum_{i=1}^{T}\lambda_i\taskvec_i$.
\For{$s=0,\ldots,S-1$}
    \State $\theta^{(s,0)} \gets \theta^{(s)}$.
    \State Choose a task order $\pi_s$ over $\{1,\ldots,T\}$.
    \For{$k=1,\ldots,T$}
        \State $i \gets \pi_s(k)$.
        \State $g_i^{(s,k)} \gets \nabla_\theta \loss_i(\theta^{(s,k-1)};\data_i^{\mathrm{probe}})$ using task $i$'s head.
        \State $v_i^{(s,k)} \gets \theta_i - \theta^{(s,k-1)}$.
        \For{each mergeable coordinate $r$}
            \State $m_{i,r}^{(s,k)} \gets \ind\{g_{i,r}^{(s,k)}v_{i,r}^{(s,k)} < -\epsilon_{\mathrm{gate}}\}$.
            \State $\widetilde u_{i,r}^{(s,k)} \gets -g_{i,r}^{(s,k)} |v_{i,r}^{(s,k)}| m_{i,r}^{(s,k)}$.
            \State $u_{i,r}^{(s,k)} \gets \clip(\eta\,\widetilde u_{i,r}^{(s,k)},- |v_{i,r}^{(s,k)}|, |v_{i,r}^{(s,k)}|)$. \Comment{clip after learning-rate scaling}
        \EndFor
        \State $\theta^{(s,k)} \gets \theta^{(s,k-1)} + u_i^{(s,k)}$.
    \EndFor
    \State $\theta^{(s+1)} \gets \theta^{(s,T)}$.
\EndFor
\State \Return $\theta^{(S)}$.
\end{algorithmic}
\end{algorithm}


\subsection{A Second Perspective: Attribution Patching in Parameter Space}
\label{sec:ap-view}
The refinement procedure developed in the previous section (Algorithm~\ref{alg:expert-gated-refinement}) can be motivated from an interpretability-style argument: it treats each move toward a task expert as a parameter-space intervention and uses a first-order attribution-patching score to decide, coordinate by coordinate, whether that intervention is locally loss-decreasing.


The attribution step is a local calculation around the current model state. To make the sequential update rule explicit, let $\vartheta$ denote the current merged encoder before applying a single task update. For task $i$, the task expert $\theta_i$ is a known reference solution, and the merge-to-expert direction is
\begin{equation}
    v_i(\vartheta) = \theta_i - \vartheta .
    \label{eq:ap-expert-direction}
\end{equation}
The corresponding parameter-space intervention is the path
\begin{equation}
    \Gamma_i(\alpha;\vartheta)
    =
    \vartheta + \alpha v_i(\vartheta),
    \qquad \alpha \in [0,1],
    \label{eq:merge-to-expert-path}
\end{equation}
where $\alpha=0$ is the current model and $\alpha=1$ is the task expert. Let
\begin{equation}
    g_i(\vartheta) = \nabla_\theta \loss_i(\vartheta;\data_i^{\mathrm{probe}})
    \label{eq:ap-gradient}
\end{equation}
be the task gradient at the current model, evaluated with task $i$'s head. A first-order attribution-patching approximation to the loss change caused by moving from $\vartheta$ toward expert $i$ is
\begin{equation}
    \Delta_i^{\mathrm{AP}}(\vartheta)
    =
    \left\langle g_i(\vartheta), v_i(\vartheta) \right\rangle
    =
    \sum_{r=1}^{d} a_{i,r}^{\mathrm{AP}}(\vartheta),
    \label{eq:ap-total}
\end{equation}
with coordinate-wise contribution
\begin{equation}
    a_{i,r}^{\mathrm{AP}}(\vartheta)
    =
    g_{i,r}(\vartheta) v_{i,r}(\vartheta).
    \label{eq:ap-coordinate}
\end{equation}
A negative contribution means that moving coordinate $r$ toward the task expert is locally predicted to reduce task $i$'s loss. A positive contribution means that the same movement is locally predicted to increase the loss. Thus the loss-decreasing expert gate is
\begin{equation}
    m_{i,r}(\vartheta)
    =
    \ind\!
    \left\{
        a_{i,r}^{\mathrm{AP}}(\vartheta) < -\epsilon_{\mathrm{gate}}
    \right\},
    \label{eq:ap-gate}
\end{equation}
where $\epsilon_{\mathrm{gate}}\geq 0$ is an optional confidence threshold; the default method uses $\epsilon_{\mathrm{gate}}=0$. The condition $g_{i,r}(\vartheta)v_{i,r}(\vartheta)<0$ is exactly the condition that the expert direction agrees with the negative gradient in coordinate $r$.

The resulting coordinate-wise update for task $i$ is a distance-scaled negative-gradient step, masked by the attribution-patching sign:
\begin{equation}
    \widetilde u_{i,r}(\vartheta)
    =
    -g_{i,r}(\vartheta)\,\left|v_{i,r}(\vartheta)\right|\,m_{i,r}(\vartheta).
    \label{eq:raw-update}
\end{equation}
Whenever $m_{i,r}(\vartheta)=1$, this update points toward the expert because $-g_{i,r}(\vartheta)$ and $v_{i,r}(\vartheta)$ have the same sign. The factor $|v_{i,r}(\vartheta)|$ makes the step small when the current model is already close to the expert in that coordinate and larger when the coordinate is far from the expert. Importantly, the gate only certifies a local first-order effect for task $i$ at the current model state. It does not by itself prove that the coordinate is harmless for other tasks, so cross-task interference must be measured experimentally.




\begin{comment}
\subsection{Finite-Difference and Mask-Stability Diagnostics}
\label{sec:diagnostic}
The method relies on a local Taylor approximation and on gradient signs estimated from very small replay buffers. The proposal should therefore validate both the attribution approximation and the stability of the induced masks. For sampled coordinates, tensors, or parameter blocks $B$, define the measured loss change
\begin{equation}
    \Delta_{i,B}^{\mathrm{meas}}(\alpha;\vartheta)
    =
    \loss_i\!\left(\vartheta + \alpha P_B v_i(\vartheta);\data_i^{\mathrm{probe}}\right)
    -
    \loss_i\!\left(\vartheta;\data_i^{\mathrm{probe}}\right),
    \label{eq:measured-delta}
\end{equation}
where $P_B$ keeps the coordinates in block $B$ and zeros the rest. The attribution-patching prediction is
\begin{equation}
    \Delta_{i,B}^{\mathrm{AP}}(\alpha;\vartheta)
    =
    \alpha \sum_{r\in B} g_{i,r}(\vartheta)v_{i,r}(\vartheta).
    \label{eq:predicted-delta}
\end{equation}
We report sign agreement between $\Delta_{i,B}^{\mathrm{AP}}$ and $\Delta_{i,B}^{\mathrm{meas}}$, Spearman correlation across sampled blocks, and precision@k for the blocks with the most negative predicted attribution. To quantify replay-buffer noise, we also compute mask stability: for two independently sampled replay buffers of size $n$, measure the agreement or Jaccard overlap between their active gates. Low stability at $n=16$ or $n=32$ would not invalidate the method, but it would suggest interpreting the gate as a noisy expert-guided regularizer rather than as a precise coordinate-level attribution.
\end{comment}


\section{Experimental Protocol}

\paragraph{Primary diagnostic setting: RoBERTa GLUE-8.}
The primary benchmark is GLUE-8: CoLA, SST-2, MRPC, STS-B, QQP, MNLI, QNLI, and RTE. We use RoBERTa-base as the default encoder. Each task expert is obtained by fine-tuning from the same pretrained checkpoint with matched hyperparameter budgets. The primary setting merges full fine-tuned encoders and keeps task-specific heads fixed for evaluation. This setting is intentionally described as \emph{oracle-task-id, multi-head encoder merging}; it is a useful and inexpensive diagnostic but not a complete demonstration of single-model deployment.

\paragraph{Recommended proof-of-concept order.}
The first experiments should be deliberately smaller than full GLUE-8, but the final report should avoid relying only on hand-picked related tasks.
\begin{enumerate}[leftmargin=*]
    \item \textbf{Code smoke test:} use a distilled RoBERTa-style encoder \citep{sanh2019distilbert} on MRPC and STS-B. This pair is small enough for rapid iteration and tests whether the implementation can exploit related sentence-pair tasks.
    \item \textbf{First convincing three-task result:} use RoBERTa-base on MRPC, STS-B, and SST-2. MRPC and STS-B form a related sentence-pair pair, while SST-2 acts as a less-related single-sentence control task.
    \item \textbf{Systematic pairwise GLUE study:} evaluate all pairwise merges among the eight GLUE tasks and report a heatmap of mean retention, worst-task retention, and expert-relative drop. This guards against cherry-picking a favorable task subset.
    \item \textbf{Structured GLUE subsets:} evaluate related subsets (MRPC/STS-B and MRPC/QQP/STS-B), a mixed-format subset (STS-B/SST-2/QNLI), small-task stress subsets (RTE/CoLA/MRPC and QQP/RTE/CoLA), and the entailment cluster (MNLI/QNLI/RTE).
    \item \textbf{Full GLUE-8 scale-up:} run the final method and all feasible baselines on all eight tasks. CoLA and RTE require multiple seeds because they are small or high-variance; MNLI requires explicit compute budgeting because it is much larger than the other tasks.
\end{enumerate}
The main paper-aligned result should use RoBERTa-base. Distilled models are useful for debugging hyperparameters and implementation errors but should not be the only evidence.

\paragraph{Secondary benchmark settings.}
To address the limitations of multi-head GLUE, we include two secondary benchmark tracks once the RoBERTa proof of concept is stable. First, a shared-output NLP track evaluates GLUE or SuperGLUE in a text-to-text format using T5-style models, so that the merged model uses a common generation head rather than task-specific classifiers \citep{wang2019superglue,raffel2020t5}. Second, a CLIP/ViT vision track evaluates the standard task-vector image-classification suite: SUN397, Cars, RESISC45, EuroSAT, SVHN, GTSRB, MNIST, and DTD \citep{radford2021clip,ilharco2022taskarithmetic,tang2024fusionbench}. The vision track is especially valuable because it tests a different modality, different loss geometry, and stronger domain shifts than GLUE. A small LoRA- or adapter-based LLM merging experiment can be added later if compute permits, but it should not replace the controlled RoBERTa and CLIP/ViT studies.

\paragraph{Replay data.}
For refinement, we use $n \in \{4,8,16,32,64,128\}$ replay examples per task. Replay examples are sampled from training data or from a held-out split separated from the final validation split. For classification tasks, replay buffers should be class-balanced whenever possible; for imbalanced tasks such as QQP and for small tasks such as RTE, we report the exact sampling procedure. We report sensitivity to the number of replay examples, variation across replay samples, and the area under the replay-budget curve rather than only the best tuned replay size.

\paragraph{Baselines.}
Because the proposed method uses labeled replay buffers, baselines are grouped by data regime.
\begin{description}[style=nextline,leftmargin=1.7em,labelindent=0pt]
    \item[Checkpoint-only baselines.] Pretrained encoder, individual single-task experts, uniform model averaging, task arithmetic, Fisher merging, TIES, DARE, DELLA-style magnitude sampling, Model Breadcrumbs, and Localize-and-Stitch if implementation allows.
    \item[Data-dependent but label-free baselines.] RegMean, AdaMerging, APL, and any FusionBench-compatible unsupervised or activation-based method that can use the same number of unlabeled examples.
    \item[Labeled-replay baselines.] Ordinary replay-buffer gradient descent from the same task-arithmetic initialization; distance-scaled gradient descent without the attribution gate; head-only replay calibration; encoder-only replay fine-tuning; encoder-plus-head replay fine-tuning; multitask replay fine-tuning from the pretrained checkpoint; and multitask replay fine-tuning from the task-arithmetic merge.
    \item[Sanity controls.] Inverted gate, random gate with matched density, wrong-expert directions, shuffled replay labels, frozen first-step gates, and the aggregated-$U$ version of the method. The aggregated-$U$ ablation is important because the proposed algorithm applies each task update immediately after it is computed.
\end{description}
All labeled-replay baselines should use the same replay examples, the same number of gradient evaluations where possible, and the same hyperparameter search budget. This is essential because the strongest alternative explanation for a positive result is simply that the method performs supervised post-merge replay fine-tuning.

\paragraph{Metrics.}
For each GLUE task, we report the standard validation metric: Matthews correlation for CoLA, accuracy for SST-2, matched/mismatched accuracy for MNLI, correlation for STS-B, and accuracy/F1-style scores for the paraphrase and question-pair tasks according to the GLUE convention. Raw task scores, expert scores, pretrained scores, task-arithmetic scores, and absolute drops from the corresponding expert are primary. Because correlation-style metrics can be near zero or negative, the normalized retention metric is
\begin{equation}
    \mathrm{NormRet}_t
    =
    \frac{
        \mathrm{Score}_t(\theta_{\mathrm{merge}})-\mathrm{Score}_t(\thetab)
    }{
        \mathrm{Score}_t(\theta_t)-\mathrm{Score}_t(\thetab)+\epsilon_{\mathrm{metric}}
    },
    \label{eq:normalized-retention}
\end{equation}
where $\epsilon_{\mathrm{metric}}$ prevents division by very small denominators. We report mean normalized retention, worst-task normalized retention, harmonic mean normalized retention, and expert-relative drop. Worst-task retention is important because a merge can improve average score while catastrophically harming a small task such as RTE or CoLA. For each main table, we also report confidence intervals across expert fine-tuning seeds and replay-buffer seeds.

\begin{comment}
\paragraph{Task-similarity and interference analysis.}
We group GLUE-8 tasks into interpretable clusters. The paraphrase and similarity cluster contains $\{\mathrm{MRPC},\mathrm{QQP},\mathrm{STS\mbox{-}B}\}$; the entailment-style cluster contains $\{\mathrm{MNLI},\mathrm{QNLI},\mathrm{RTE}\}$; and the single-sentence cluster contains $\{\mathrm{CoLA},\mathrm{SST\mbox{-}2}\}$. We test whether expert-direction-gated refinement helps most in related task clusters, where movement toward one expert is more likely to be compatible with other tasks. We also compute task-vector cosine similarity, cross-task transfer performance, average attribution-patching sign agreement, and the cross-task interference matrix
\begin{equation}
    C_{j\leftarrow i}^{(s,k)}
    =
    \left\langle
        \nabla_\theta \loss_j(\theta^{(s,k-1)};\data_j^{\mathrm{probe}}),
        u_i^{(s,k)}
    \right\rangle.
    \label{eq:interference-matrix}
\end{equation}
Negative entries indicate that task $i$'s update is locally predicted to help task $j$, while positive entries indicate predicted interference. This matrix directly tests whether the gate is only self-protective or also tends to align across related tasks.
\end{comment}

\subsection{Ablations and Diagnostics}
We plan the following ablations:
\begin{enumerate}[leftmargin=*]
    \item \textbf{Gate sign and necessity:} compare the proposed gate $g_{i,r}v_{i,r}<0$ with no gate, an inverted gate $g_{i,r}v_{i,r}>0$, and a random mask of the same density.
    \item \textbf{Gate confidence and granularity:} vary $\epsilon_{\mathrm{gate}}$, compare coordinate-wise gates with tensor-wise or block-wise gates, evaluate top-$k$ most negative attribution masks, and test sign-stability masks formed from multiple replay minibatches.
    \item \textbf{Sequential update rule:} compare the proposed immediate single-task update with the aggregated-$U$ variant, fixed task order, cyclic task order, and random task order.
    \item \textbf{Distance scaling:} compare ordinary replay-gradient descent, distance-scaled gradient descent without the gate, direct expert-interpolation steps $u_i\propto m_i v_i$, and the full gated update.
    \item \textbf{Clipping and normalization:} compare no normalization, tensor-wise RMS normalization, per-task global norm clipping, and post-update validation line search.
    \item \textbf{Replay budget and variance:} vary the number of replay examples per task and report variation across replay samples and random seeds.
    \item \textbf{Number of refinement steps:} compare $S\in\{0,1,2,5,10\}$ to test whether most gains occur in the first few steps or whether overfitting appears.
    \item \textbf{Expansion point:} recompute gradients and gates before every single-task update versus freezing the first-sweep gates from $\theta^{(0)}$.
    \item \textbf{Label and expert sanity checks:} compare correct labels with shuffled replay labels, and correct expert directions with wrong-expert directions.
\end{enumerate}

\begin{comment}
\subsection{Hypotheses}
The project evaluates five hypotheses:
\begin{description}[style=nextline,leftmargin=1.5em,labelindent=0pt]
    \item[H1: Attribution-patching signs predict useful expert directions.] Coordinates or blocks with $g_{i,r}v_{i,r}<0$ should be more likely to reduce task $i$'s loss when moved toward expert $i$ than coordinates or blocks with $g_{i,r}v_{i,r}>0$.
    \item[H2: Expert-guided replay refinement improves task arithmetic.] Starting from the same task-arithmetic initialization and using the same replay budget, the gated update should outperform ordinary replay-buffer gradient descent and distance-scaled gradient descent without the gate.
    \item[H3: Sequential single-task updates improve worst-task retention.] Applying each clipped task update immediately after it is computed should reduce overshoot and stale-gradient effects relative to an aggregated-$U$ update, improving worst-task retention under the same replay budget.
    \item[H4: Benefits are larger for related tasks.] The method should help most in related task clusters, especially MRPC, QQP, and STS-B, where expert directions are more likely to be mutually compatible.
    \item[H5: Gains are not GLUE-specific.] A positive GLUE result should transfer at least partially to a shared-head text-to-text setting or to the CLIP/ViT vision task-vector suite.
\end{description}
\end{comment}

\section{Preliminary Results}
\label{sec:results}
We report initial experiments in the RoBERTa-base multi-head setting (oracle task identity at evaluation). Unless noted, refinement uses $S$ sweeps from the task-arithmetic merge with the corrected trust region applied \emph{after} the learning rate (Eq.~\eqref{eq:clipped-update}) and threshold $\epsilon_{\mathrm{gate}}=0$. Each method family (ordinary replay gradient descent, the proposed gated update, and the ungated distance-scaled update) is given its own learning-rate search of equal size; we report the best configuration per family by mean normalized retention (Eq.~\eqref{eq:normalized-retention}). These results are preliminary: the three-task study uses five seeds, but the GLUE-8 study is currently single-seed, $\lambda_i=0.3$ is untuned, and the multi-head setting assumes oracle task identity.



\paragraph{Three-task proof of concept.}
On the MRPC/STS-B/SST-2 triple ($n=64$ replay examples per task, $S=5$), all refinement methods recover most of the gap from the task-arithmetic merge (mean normalized retention $0.558$) toward the experts, and the three families are statistically tied at their tuned optima (Table~\ref{tab:poc3}). The decisive difference is robustness to the learning rate: the gated update attains its best score at a much larger step ($\eta=64$) and remains stable there, whereas the ungated update---identical except for the gate---\emph{diverges} at the same learning rate (mean normalized retention $-1.07$). The attribution gate thus acts as a protective mask that enables large, safe steps by suppressing the locally loss-increasing expert directions.

\begin{table}[t]
\centering
\caption{Three-task RoBERTa-base proof of concept (MRPC, STS-B, SST-2; $n=64$, $S=5$). Mean normalized retention; best learning rate per family. Left: best configuration, mean $\pm$ std over five seeds. Right: a single-seed contrast at a large fixed learning rate, isolating the gate's effect.}
\label{tab:poc3}
\begin{tabular}{lcc}
\toprule
Method & Best mean NormRet ($\pm$ std) & at $\eta=64$ (1 seed) \\
\midrule
Task-arithmetic merge & $0.558$ & --- \\
Ordinary replay GD & $0.831\pm0.005$ & --- \\
Ungated distance-scaled (no gate) & $0.835\pm0.017$ & $-1.07$ (diverges) \\
\textbf{APR (gated, proposed)} & $\mathbf{0.840\pm0.021}$ & $\mathbf{0.849}$ (stable) \\
\bottomrule
\end{tabular}
\end{table}

\paragraph{GLUE-8 scale-up.}
Merging all eight experts (one consistent RoBERTa-base fine-tuning source) is a much harder, more interference-prone setting: the task-arithmetic merge attains only mean normalized retention $0.320$, with MRPC driven below the pretrained baseline. Here the value of expert anchoring grows: both anchored updates clearly outperform ordinary replay gradient descent, by a larger margin than in the three-task case and especially on worst-task retention (Table~\ref{tab:glue8}). The gate is approximately performance-neutral relative to the ungated anchored update on mean retention---the ordering flips between the two budgets and the gaps are within single-seed noise---but it again provides a wider safe learning-rate range (the gated update peaks at $\eta=32$ and degrades, while the ungated update peaks at $\eta=2\text{--}4$ and diverges by $\eta\gtrsim 8$). MRPC and CoLA remain the persistent worst tasks; QQP, STS-B, MNLI, and QNLI recover to $0.78\text{--}0.86$.

\begin{table}[t]
\centering
\caption{GLUE-8 (all eight tasks, RoBERTa-base, single seed). Best mean normalized retention per family (worst-task retention in parentheses); the task-arithmetic merge baseline is $0.320$ (worst $-0.175$).}
\label{tab:glue8}
\begin{tabular}{lcc}
\toprule
Method & $n{=}64,\,S{=}5$ & $n{=}32,\,S{=}10$ \\
\midrule
Ordinary replay GD & $0.574\ (0.04)$ & $0.588\ (0.04)$ \\
Ungated distance-scaled (no gate) & $\mathbf{0.637}\ (0.15)$ & $0.607\ (0.05)$ \\
\textbf{APR (gated, proposed)} & $0.614\ (0.14)$ & $\mathbf{0.617}\ (0.05)$ \\
\bottomrule
\end{tabular}
\end{table}

\paragraph{What the attribution gate is and is not doing.}
Two diagnostics temper the interpretation of the gate. First, the gate keeps a strikingly stable $\approx 36\%$ of coordinates across sweeps, tasks, and learning rates. A control attributes most of this to gradient sparsity rather than to the expert geometry: roughly $31\%$ of encoder coordinates are word-embedding rows for tokens absent from the small replay buffer and therefore have exactly zero gradient (always masked), and a random direction of matched magnitudes reproduces almost the same density ($0.345$ vs.\ $0.361$ for the true expert direction). The expert direction itself contributes only a small signal beyond chance. Second, the trust-region clip is almost never active at useful learning rates ($<0.1\%$ of gated coordinates), functioning as a rare-event safeguard against the few large-gradient coordinates rather than as a routine operation. The gate's measurable benefit is therefore concentrated in learning-rate robustness---masking the loss-increasing active coordinates so that large steps remain safe---rather than in a higher peak score at a tuned learning rate.

\paragraph{CLIP/ViT vision track.}
We next test whether the method transfers beyond GLUE encoders to the standard eight-task CLIP task-vector suite (SUN397, Cars, RESISC45, EuroSAT, SVHN, GTSRB, MNIST, DTD) with a CLIP ViT-B/32 backbone. Here the mergeable model is the image encoder, and each task keeps a \emph{frozen} zero-shot classification head derived from the CLIP text encoder over class-name prompts; the experts are publicly available fine-tuned image encoders, merged with uniform $\lambda_i=0.3$. Refinement uses $n=64$, $S=5$, and the trust region of Eq.~\eqref{eq:clipped-update} with $\gamma=1$ (single seed). This is a stronger test than GLUE: a different modality, dense gradients without the word-embedding sparsity discussed above, and larger domain shifts. The task-arithmetic merge is weak in this regime (mean normalized retention $0.270$, worst-task $-0.527$; SUN397 and Cars are pushed \emph{below} the zero-shot backbone). All three families recover substantial ground, but---in contrast to GLUE-8---the attribution gate now attains both the best mean retention and the only positive worst-task retention (Table~\ref{tab:clip8}). The gated update improves the merge from $0.270$ to $0.598$ mean and from $-0.527$ to $+0.159$ worst-task, beating ordinary replay gradient descent by $+0.14$ mean / $+0.36$ worst and the ungated anchored update by $+0.06$ mean / $+0.25$ worst. The gains concentrate on the most interference-damaged tasks: Cars and SUN397, which the merge drove below the zero-shot floor, are precisely the ones the gate rescues (Table~\ref{tab:clip8pt}), whereas the ungated update sacrifices Cars catastrophically at its larger steps (worst-task $-2.97$ at $\eta=8$). An inverted-gate control diverges (mean $-1.68$), confirming that the sign of the attribution product carries the signal. A smaller three-task vision study (EuroSAT/GTSRB/MNIST, single seed) shows the same ordering more sharply: APR $0.940$ mean / $0.915$ worst, versus $0.917/0.876$ for the ungated update, $0.874/0.792$ for ordinary replay GD, and $0.811/0.721$ for the merge. As in GLUE, the best learning rate falls as tasks are added (the gated update peaks at $\eta=32$ on three tasks and at $\eta=8$ on eight). Together these results support H5---the method transfers to a second modality---and indicate that the gate's usefulness is \emph{modality-dependent}: approximately neutral on GLUE encoders but beneficial, on both mean and worst-task retention, on CLIP/ViT.

\begin{table}[t]
\centering
\caption{Eight-task CLIP ViT-B/32 suite (single seed, $n=64$, $S=5$, $\gamma=1$, uniform $\lambda_i=0.3$). Best learning rate per family by mean normalized retention (Eq.~\eqref{eq:normalized-retention}); worst-task retention in parentheses. The gated update is best on both, and is the only method keeping worst-task retention positive.}
\label{tab:clip8}
\begin{tabular}{lcc}
\toprule
Method & Mean NormRet & Worst-task \\
\midrule
Task-arithmetic merge & $0.270$ & $-0.527$ \\
Ordinary replay GD & $0.458$ & $-0.203$ \\
Ungated distance-scaled (no gate) & $0.536$ & $-0.086$ \\
\textbf{APR (gated, proposed)} & $\mathbf{0.598}$ & $\mathbf{+0.159}$ \\
\midrule
Inverted-gate control & $-1.682$ & $-5.386$ \\
\bottomrule
\end{tabular}
\end{table}

\begin{table}[t]
\centering
\footnotesize
\setlength{\tabcolsep}{4pt}
\caption{Per-task top-1 accuracy on the eight-task CLIP suite: zero-shot backbone, task expert, task-arithmetic merge, and the proposed gated refinement (APR, $\eta=8$). Refinement recovers most of the merge's interference loss, including on Cars and SUN397, which the merge had degraded below the zero-shot backbone.}
\label{tab:clip8pt}
\begin{tabular}{lccccccccc}
\toprule
 & SUN397 & Cars & RESISC & EuroSAT & SVHN & GTSRB & MNIST & DTD & Avg. \\
\midrule
Zero-shot & $0.632$ & $0.597$ & $0.582$ & $0.450$ & $0.315$ & $0.325$ & $0.482$ & $0.439$ & $0.478$ \\
Expert    & $0.749$ & $0.783$ & $0.949$ & $0.991$ & $0.972$ & $0.989$ & $0.996$ & $0.794$ & $0.903$ \\
Merge     & $0.570$ & $0.553$ & $0.628$ & $0.734$ & $0.778$ & $0.685$ & $0.961$ & $0.472$ & $0.673$ \\
\textbf{APR} & $\mathbf{0.650}$ & $\mathbf{0.648}$ & $\mathbf{0.769}$ & $\mathbf{0.923}$ & $\mathbf{0.874}$ & $\mathbf{0.851}$ & $0.961$ & $\mathbf{0.578}$ & $\mathbf{0.782}$ \\
\bottomrule
\end{tabular}
\end{table}

\paragraph{Replay budget and refinement horizon (CLIP/ViT).}
Two further studies on the eight-task suite probe the labeled-replay budget and the number of refinement sweeps, comparing the gated update directly against ordinary replay gradient descent---the strongest baseline that also uses labeled replay. First, reducing the replay buffer from $n=64$ to $32$ and $16$ examples per task degrades both methods, but the gated update retains an almost constant advantage of ${\approx}\,+0.10$ mean normalized retention at every budget and remains the only method with non-negative worst-task retention down to $n=16$ (Table~\ref{tab:clipbudget}). Its best learning rate also decreases as the buffer shrinks (from $\eta=8$ at $n=64$ to $\eta=6$ at $n\le 32$), consistent with noisier replay gradients favoring smaller trust-region steps. Second, on the refinement horizon: the gated update is essentially converged after five sweeps (mean retention $0.598\to0.615$ from $S=5$ to $S=15$ at $n=64$, with no sign of overfitting the buffer), whereas ordinary gradient descent is still improving at $S=35$ ($0.458\to0.504\to0.526$ at $S=5,15,35$) yet does not catch up and---critically---never lifts its worst task (SUN397) back above the zero-shot backbone. Thus the expert anchor, rather than additional replay data or gradient steps, is what recovers the most interference-damaged tasks.

\begin{table}[t]
\centering
\caption{Replay-budget study on the eight-task CLIP suite (mean normalized retention; worst-task in parentheses). Best configuration per family; the gated update (APR) near-converges by $S=5\text{--}10$ sweeps, while ordinary replay GD is given a longer horizon ($S=25\text{--}35$) as a generous baseline. APR holds an ${\approx}\,+0.10$ mean advantage at every budget and is the only method with non-negative worst-task retention down to $n=16$.}
\label{tab:clipbudget}
\begin{tabular}{lcc}
\toprule
Replay budget & \textbf{APR (ours)} & Ordinary replay GD \\
\midrule
$n=16$ & $\mathbf{0.522}\ (\mathbf{+0.018})$ & $0.427\ (-0.266)$ \\
$n=32$ & $\mathbf{0.570}\ (\mathbf{+0.122})$ & $0.476\ (-0.197)$ \\
$n=64$ & $\mathbf{0.598}\ (\mathbf{+0.159})$ & $0.526\ (-0.135)$ \\
\bottomrule
\end{tabular}
\end{table}



\begin{thebibliography}{99}

\bibitem[Wang et~al.(2018)]{wang2018glue}
Alex Wang, Amanpreet Singh, Julian Michael, Felix Hill, Omer Levy, and Samuel R. Bowman.
\newblock GLUE: A multi-task benchmark and analysis platform for natural language understanding.
\newblock In \emph{Proceedings of the 2018 EMNLP Workshop BlackboxNLP}, 2018.

\bibitem[Liu et~al.(2019)]{liu2019roberta}
Yinhan Liu, Myle Ott, Naman Goyal, Jingfei Du, Mandar Joshi, Danqi Chen, Omer Levy, Mike Lewis, Luke Zettlemoyer, and Veselin Stoyanov.
\newblock RoBERTa: A robustly optimized BERT pretraining approach.
\newblock \emph{arXiv preprint arXiv:1907.11692}, 2019.

\bibitem[Wang et~al.(2019)]{wang2019superglue}
Alex Wang, Yada Pruksachatkun, Nikita Nangia, Amanpreet Singh, Julian Michael, Felix Hill, Omer Levy, and Samuel R. Bowman.
\newblock SuperGLUE: A stickier benchmark for general-purpose language understanding systems.
\newblock In \emph{Advances in Neural Information Processing Systems}, 2019.

\bibitem[Raffel et~al.(2020)]{raffel2020t5}
Colin Raffel, Noam Shazeer, Adam Roberts, Katherine Lee, Sharan Narang, Michael Matena, Yanqi Zhou, Wei Li, and Peter J. Liu.
\newblock Exploring the limits of transfer learning with a unified text-to-text transformer.
\newblock \emph{Journal of Machine Learning Research}, 21(140):1--67, 2020.

\bibitem[Radford et~al.(2021)]{radford2021clip}
Alec Radford, Jong Wook Kim, Chris Hallacy, Aditya Ramesh, Gabriel Goh, Sandhini Agarwal, Girish Sastry, Amanda Askell, Pamela Mishkin, Jack Clark, Gretchen Krueger, and Ilya Sutskever.
\newblock Learning transferable visual models from natural language supervision.
\newblock In \emph{International Conference on Machine Learning}, 2021.


\bibitem[Sanh et~al.(2019)]{sanh2019distilbert}
Victor Sanh, Lysandre Debut, Julien Chaumond, and Thomas Wolf.
\newblock DistilBERT, a distilled version of BERT: smaller, faster, cheaper and lighter.
\newblock \emph{arXiv preprint arXiv:1910.01108}, 2019.

\bibitem[Ilharco et~al.(2022)]{ilharco2022taskarithmetic}
Gabriel Ilharco, Marco Tulio Ribeiro, Mitchell Wortsman, Suchin Gururangan, Ludwig Schmidt, Hannaneh Hajishirzi, and Ali Farhadi.
\newblock Editing models with task arithmetic.
\newblock \emph{arXiv preprint arXiv:2212.04089}, 2022.

\bibitem[Wortsman et~al.(2022)]{wortsman2022modelsoups}
Mitchell Wortsman, Gabriel Ilharco, Samir Yitzhak Gadre, Rebecca Roelofs, Raphael Gontijo-Lopes, Ari S. Morcos, Hongseok Namkoong, Ali Farhadi, Yair Carmon, Simon Kornblith, and Ludwig Schmidt.
\newblock Model soups: Averaging weights of multiple fine-tuned models improves accuracy without increasing inference time.
\newblock In \emph{International Conference on Machine Learning}, 2022.

\bibitem[Matena and Raffel(2022)]{matena2022fisher}
Michael S. Matena and Colin A. Raffel.
\newblock Merging models with Fisher-weighted averaging.
\newblock In \emph{Advances in Neural Information Processing Systems}, 2022.

\bibitem[Jin et~al.(2023)]{jin2023regmean}
Xisen Jin, Xiang Ren, Daniel Preotiuc-Pietro, and Pengxiang Cheng.
\newblock Dataless knowledge fusion by merging weights of language models.
\newblock In \emph{International Conference on Learning Representations}, 2023.

\bibitem[Yadav et~al.(2023)]{yadav2023ties}
Prateek Yadav, Derek Tam, Leshem Choshen, Colin Raffel, and Mohit Bansal.
\newblock TIES-Merging: Resolving interference when merging models.
\newblock \emph{arXiv preprint arXiv:2306.01708}, 2023.

\bibitem[Yu et~al.(2023)]{yu2023dare}
Le Yu, Bowen Yu, Haiyang Yu, Fei Huang, and Yongbin Li.
\newblock Language models are Super Mario: Absorbing abilities from homologous models as a free lunch.
\newblock \emph{arXiv preprint arXiv:2311.03099}, 2023.

\bibitem[Davari and Belilovsky(2023)]{davari2023breadcrumbs}
Mohammad Davari and Eugene Belilovsky.
\newblock Model Breadcrumbs: Scaling multi-task model merging with sparse masks.
\newblock \emph{arXiv preprint arXiv:2312.06795}, 2023.

\bibitem[Yang et~al.(2024)]{yang2024adamerging}
Enneng Yang, Zhenyi Wang, Li Shen, Shiwei Liu, Guibing Guo, Xingwei Wang, and Dacheng Tao.
\newblock AdaMerging: Adaptive model merging for multi-task learning.
\newblock \emph{arXiv preprint arXiv:2310.02575}, 2024.

\bibitem[Deep et~al.(2024)]{deep2024della}
Pala Tej Deep, Rishabh Bhardwaj, and Soujanya Poria.
\newblock DELLA-Merging: Reducing interference in model merging through magnitude-based sampling.
\newblock \emph{arXiv preprint arXiv:2406.11617}, 2024.

\bibitem[Kong et~al.(2024)]{kong2024apl}
Fanshuang Kong, Richong Zhang, and Ziqiao Wang.
\newblock Activated parameter locating via causal intervention for model merging.
\newblock \emph{arXiv preprint arXiv:2408.09485}, 2024.

\bibitem[He et~al.(2024)]{he2024localize}
Shwai He, Runqi Wang, Jindong Wang, Jindong Gu, and Xing Xie.
\newblock Localize-and-Stitch: Efficient model merging via sparse task arithmetic.
\newblock \emph{arXiv preprint arXiv:2408.13656}, 2024.

\bibitem[Lu et~al.(2024)]{lu2024twin}
Zhenyi Lu, Chengxu Yin, Wujie Wang, and Xuebo Liu.
\newblock Twin-Merging: Dynamic integration of modular expertise in model merging.
\newblock \emph{arXiv preprint arXiv:2406.15479}, 2024.

\bibitem[Tang et~al.(2024)]{tang2024fusionbench}
Anke Tang, Li Shen, Yong Luo, Liang Ding, Han Hu, Bo Du, and Dacheng Tao.
\newblock FusionBench: A comprehensive benchmark of deep model fusion.
\newblock \emph{arXiv preprint arXiv:2406.03280}, 2024.

\bibitem[Zeng et~al.(2025)]{zeng2025taskvectorbases}
Zihan Zeng, Jiahui Chen, Lei Li, and Min Zhang.
\newblock Efficient model editing with task vector bases: A theoretical framework and scalable approach.
\newblock \emph{arXiv preprint arXiv:2502.01015}, 2025.

\bibitem[Sundararajan et~al.(2017)]{sundararajan2017ig}
Mukund Sundararajan, Ankur Taly, and Qiqi Yan.
\newblock Axiomatic attribution for deep networks.
\newblock In \emph{International Conference on Machine Learning}, 2017.

\bibitem[Syed et~al.(2024)]{syed2024attributionpatching}
Aaquib Syed, Can Rager, and Arthur Conmy.
\newblock Attribution patching outperforms automated circuit discovery.
\newblock \emph{arXiv preprint arXiv:2310.10348}, 2024.

\bibitem[Kramar et~al.(2024)]{kramar2024atpstar}
Janos Kramar, Tom Lieberum, Rohin Shah, and Neel Nanda.
\newblock AtP*: An efficient and scalable method for localizing language model behavior to components.
\newblock \emph{arXiv preprint arXiv:2403.00745}, 2024.

\end{thebibliography}










\section{Appendix}

\begin{theorem}[Box-distance monotonicity under coordinatewise movement toward a hull point]
\label{thm:box-distance-monotonicity}
Let \(x_1,\dots,x_n\in \mathbb{R}^D\). Define the axis-aligned box hull
\[
B:=\prod_{d=1}^D [m_d,M_d],
\qquad
m_d:=\min_{1\le j\le n} x_j^d,
\qquad
M_d:=\max_{1\le j\le n} x_j^d.
\]
Let \(y,z\in \mathbb{R}^D\). Suppose that there exists some
\(i\in\{1,\dots,n\}\) such that, for every coordinate \(d\in\{1,\dots,D\}\),
the point \(z^d\) lies between \(x_i^d\) and \(y^d\). Equivalently,
\[
|z^d-x_i^d|\le |y^d-x_i^d|
\]
and \(z^d-x_i^d\) has the same sign as \(y^d-x_i^d\), allowing equality.

Then, for every \(1\le p\le \infty\),
\[
d_p(z,B)\le d_p(y,B),
\]
where \(d_p(u,B)\) denotes the distance from \(u\) to \(B\) with respect to
the \(\ell^p\)-norm.
\end{theorem}

\begin{proof}
For each coordinate \(d\), set
\[
I_d:=[m_d,M_d].
\]
Since \(x_i\in B\), we have
\[
x_i^d\in I_d
\]
for every \(d\).

Define the one-dimensional distance function
\[
\delta_d(t):=\operatorname{dist}(t,I_d)
=
\inf_{s\in I_d}|t-s|.
\]

We first prove the following elementary one-dimensional claim.

\medskip

\noindent
\textbf{Claim.}
Let \(I=[m,M]\subset \mathbb{R}\), let \(c\in I\), and suppose that
\(z\) lies between \(c\) and \(y\). Then
\[
\operatorname{dist}(z,I)\le \operatorname{dist}(y,I).
\]

\medskip

\noindent
Indeed, if \(y\in I\), then since \(c\in I\) and \(I\) is an interval,
every point between \(c\) and \(y\) also lies in \(I\). Hence \(z\in I\), so
\[
\operatorname{dist}(z,I)=0=\operatorname{dist}(y,I).
\]

If \(y>M\), then
\[
\operatorname{dist}(y,I)=y-M.
\]
Since \(z\) lies between \(c\in I\) and \(y>M\), we have \(z\le y\). If
\(z\le M\), then \(\operatorname{dist}(z,I)=0\). If \(z>M\), then
\[
\operatorname{dist}(z,I)=z-M\le y-M=\operatorname{dist}(y,I).
\]
The case \(y<m\) is symmetric. This proves the claim.

\medskip

Now apply the claim coordinatewise with
\[
I=I_d,
\qquad
c=x_i^d.
\]
Because \(z^d\) lies between \(x_i^d\) and \(y^d\), we obtain
\[
\delta_d(z^d)\le \delta_d(y^d)
\]
for every \(d\in\{1,\dots,D\}\).

Next, distance to an axis-aligned box separates coordinatewise. Namely, for
\(1\le p<\infty\),
\[
d_p(u,B)
=
\left(
\sum_{d=1}^D \delta_d(u^d)^p
\right)^{1/p},
\]
and for \(p=\infty\),
\[
d_\infty(u,B)
=
\max_{1\le d\le D}\delta_d(u^d).
\]
This follows because the closest point of \(B\) to \(u\) is obtained by
clipping each coordinate \(u^d\) to the interval \(I_d=[m_d,M_d]\).

Since
\[
\delta_d(z^d)\le \delta_d(y^d)
\]
for every \(d\), monotonicity of the \(\ell^p\)-norm on the nonnegative
orthant gives
\[
d_p(z,B)\le d_p(y,B).
\]
This proves the theorem.
\end{proof}




\begin{proposition}
\label{prop:l1-geodesic-box-hull}
The set \(B\) is the \(\ell^1\)-geodesic convex hull of
\(\{x_1,\dots,x_n\}\).
\end{proposition}

\begin{proof}
For \(a,b\in \mathbb{R}^D\), define the \(\ell^1\)-metric interval between
\(a\) and \(b\) by
\[
I_1(a,b)
:=
\{u\in\mathbb{R}^D:
\|a-u\|_1+\|u-b\|_1=\|a-b\|_1\}.
\]
We claim that
\[
I_1(a,b)
=
\prod_{d=1}^D
[\min(a^d,b^d),\max(a^d,b^d)].
\]

Indeed, by the triangle inequality in one dimension,
\[
|a^d-u^d|+|u^d-b^d|\ge |a^d-b^d|
\]
for each coordinate \(d\). Summing over \(d\), equality in the
\(\ell^1\)-triangle inequality holds if and only if equality holds in every
coordinate. In one dimension, equality holds precisely when \(u^d\) lies
between \(a^d\) and \(b^d\). Hence the displayed formula for \(I_1(a,b)\)
follows.

Now \(B\) is clearly \(\ell^1\)-geodesically convex: if \(a,b\in B\), then
every point coordinatewise between \(a\) and \(b\) also lies in \(B\).

It remains to show minimality. Let \(C\subseteq \mathbb{R}^D\) be any
\(\ell^1\)-geodesically convex set containing \(x_1,\dots,x_n\). Since
\(C\) is closed under \(\ell^1\)-metric intervals, it is closed under taking
coordinatewise boxes between pairs of its points. In particular, if
\(a,b\in C\), then both the coordinatewise minimum
\[
a\wedge b
:=
(\min(a^1,b^1),\dots,\min(a^D,b^D))
\]
and the coordinatewise maximum
\[
a\vee b
:=
(\max(a^1,b^1),\dots,\max(a^D,b^D))
\]
belong to \(C\).

By induction, \(C\) therefore contains the coordinatewise minimum
\[
m:=(m_1,\dots,m_D)
\]
and the coordinatewise maximum
\[
M:=(M_1,\dots,M_D)
\]
of the points \(x_1,\dots,x_n\). Since \(C\) is \(\ell^1\)-geodesically
convex, it contains the entire \(\ell^1\)-metric interval between \(m\) and
\(M\). But
\[
I_1(m,M)
=
\prod_{d=1}^D [m_d,M_d]
=
B.
\]
Therefore every \(\ell^1\)-geodesically convex set containing
\(x_1,\dots,x_n\) also contains \(B\). Hence \(B\) is the smallest such set,
i.e.
\[
B=\operatorname{gconv}_{\ell^1}\{x_1,\dots,x_n\}.
\]
\end{proof}









\end{document}